\documentclass[10pt,a4paper]{amsart}
\usepackage[margin=19mm]{geometry}
\usepackage{amsmath,amssymb,mathtools}
\usepackage[scr=rsfs,cal=euler]{mathalfa}
\usepackage{xfrac}
\usepackage[unicode,hidelinks,bookmarks=false]{hyperref}
\newcommand{\ie}{i.e.\ }
\newcommand{\cf}{cf.\ }
\newcommand{\resp}{respectively\ }
\newcommand{\Subsection}{\subsection}
\providecommand{\nobreakdash}{\nobreak\hbox{-}}
\setlength{\emergencystretch}{2em}
\multlinegap0pt
\usepackage{bm}
\usepackage{bbm}
\usepackage{dsfont}
\usepackage{xspace}
\usepackage{cancel}
\usepackage{mathtools}
\usepackage{amsthm}
\makeatletter
\@ifundefined{theorem}{\newtheorem{theorem}{Theorem}}{}
\@ifundefined{definition}{\newtheorem{definition}[theorem]{Definition}}{}
\@ifundefined{remark}{\newtheorem{remark}[theorem]{Remark}}{}
\makeatother
\newcommand{\C}{\ensuremath{\mathbb{C}}}
\newcommand{\R}{\ensuremath{\mathbb{R}}}
\newcommand{\T}{\ensuremath{\mathbb{T}}}
\newcommand{\Z}{\ensuremath{\mathbb{Z}}}
\newcommand{\p}{\ensuremath{\mathbb{P}}}
\title[Equilibrium stability for a continuous time scale with discr]{Equilibrium stability for a continuous time scale with discrete uniform gaps}
\author{Douglas R. Anderson}
\date{Source: arXiv:2606.02635v1 (CC BY 4.0)}
\begin{document}
\maketitle

\noindent\emph{Source and licence.} Equilibrium stability for a continuous time scale with discrete uniform gaps. Version arXiv:2606.02635v1 (CC BY 4.0), distributed under the Creative Commons Attribution 4.0 International licence, as recorded in the arXiv licence metadata for this exact version and shown on the abstract page https://arxiv.org/abs/2606.02635v1. Full rights notice: licenses/arxiv-2606.02635-CC-BY-4.0.txt.\par\smallskip

\noindent\textit{Editorial scope.} Counted unit: the Theorem whose source label is \texttt{complexTHMHUS} in the article (source0.tex lines 468--482, source label \texttt{complexTHMHUS}), delivered together with the article's own complete original proof of it (source0.tex lines 484--561, one \texttt{proof} environment of 78 lines). No mathematics is written, repaired or re-derived: statement and proof are the article's own text line for line. Required context retained uncounted: maineq (lines 116--122); pabexp0 (lines 131--133); xform (lines 135--137); reallambdathm (lines 159--187); nablaexpo (lines 255--267); nabeq (lines 263--265); complexlam (lines 327--329); complexTHM (lines 336--343); complexnab (lines 380--390); phiineq (lines 412--418). Every definition, display and label of those lines is retained unchanged. Omitted from this package and not redistributed: 1--115, 123--130, 134--134, 138--158, 188--254, 266--326, 330--335, 344--379, 391--411, 419--467, 483--483, 562--651. Nothing inside a retained excerpt is omitted or altered: every retained body is delivered exactly as the article's lines are written, so no omission receipt of any kind accompanies this package. Citations: the retained excerpts carry no citation command at all, so no bibliography entry of the article is transcribed and no reference is redistributed. Reference adaptation: none was needed and none was made; every reference inside the delivered unit resolves inside it, so no unresolved reference or citation is produced. Printed slips kept literal: no printed word, symbol, hypothesis or number of the counted statement or of its proof was altered. Layout and class: the article's own class is \texttt{amsart} with packages including amssymb, amsmath, amsfonts, amsthm, mathrsfs, graphicx, fontenc; the wrapper is the lane's pinned \texttt{amsart} editorial wrapper at 10pt on A4, which restates the theorem-like environments the retained text needs and adds the article's own macro definitions that the retained text needs (\texttt{\textbackslash C}, \texttt{\textbackslash R}, \texttt{\textbackslash T}, \texttt{\textbackslash Z}, \texttt{\textbackslash p}), with extra wrapper packages (bm, bbm, dsfont, xspace, cancel, mathtools). The delivered numbering is the unit's own. Assets: no retained span contains a figure environment, a graphics-inclusion command, a PDF-page inclusion, a table or a TikZ picture; no image file is copied into this package, the package declares no asset dependency and no image file is redistributed with it. Licence: version arXiv:2606.02635v1 of this article is distributed under the Creative Commons Attribution 4.0 International licence, as recorded in the arXiv licence metadata for this exact version and shown on the abstract page https://arxiv.org/abs/2606.02635v1. A full rights notice is delivered at licenses/arxiv-2606.02635-CC-BY-4.0.txt. Characters: every retained excerpt is pure ASCII, so the delivered engine, XeLaTeX with its default Latin Modern text font, reports no missing character.
\begin{definition}[Equilibrium Stability]
Consider the dynamic equation given by
\begin{equation}\label{maineq}
 x^\Delta(t) = \lambda x(t), \quad \lambda\in\C\backslash\left\{-\frac{1}{\beta}\right\}, \quad t\in\p_{\alpha,\beta}.
\end{equation}
We say that the equilibrium (trivial) solution of \eqref{maineq} is stable on $\p_{\alpha,\beta}$ if and only if given $\varepsilon>0$, there exists a $\delta>0$ such that if $x$ is a solution of \eqref{maineq} with $|x(0)|=|x_0|<\delta$, then for all $t\in\p_{\alpha,\beta}$ we have $|x(t)|<\varepsilon$ on $\p_{\alpha,\beta}$. Additionally, we say the equilibrium solution of \eqref{maineq} is asymptotically stable on $\p_{\alpha,\beta}$ if and only if there exists $\delta>0$ such that if $|x_0|<\delta$, then $\displaystyle\lim_{t\rightarrow\infty} |x(t)|=0$. The equilibrium solution of \eqref{maineq} is unstable if it is not stable.
\end{definition}

\begin{equation}\label{pabexp0} 
 e_{\lambda}(t,0)=\left[(1+\beta\lambda)e^{\alpha\lambda}\right]^{k}e^{\lambda j}, \quad t=k(\alpha+\beta)+j, \quad j\in[0,\alpha]. 
\end{equation}

\begin{equation}\label{xform}
 x(t)=x_0e_{\lambda}(t,0), \quad t\in\T,
\end{equation}

\begin{theorem}[Delta equation]\label{reallambdathm}
Fix $\alpha>0$, and let $\lambda\in\R\backslash\left\{-\frac{1}{\beta}\right\}$. We have the following cases.
\begin{enumerate}
\item Consider a small gap size $\beta$ that satisfies $0<\beta<\frac{\alpha}{W_0(e^{-1})}$.
 \begin{enumerate}
   \item If $\lambda\in\left(-\infty,\; -\frac{1}{\beta}\right)\cup\left(-\frac{1}{\beta},0\right)$, then the trivial solution of \eqref{maineq} is asymptotically stable. 
   \item If $\lambda=0$, then \eqref{maineq} is stable. 
   \item If $\lambda\in\left(0,\infty\right)$, then \eqref{maineq} is unstable.
 \end{enumerate}
\item Consider a critical gap size $\beta$ that satisfies $\beta=\frac{\alpha}{W_0(e^{-1})}$. Then, $(1+\beta\lambda)e^{\alpha\lambda}=-1$ at $\lambda=-\frac{1}{\alpha}\left(1+W_0(e^{-1})\right)$, and we have the following subcases.
 \begin{enumerate}
   \item If $\lambda\in\left(-\infty,\; -\frac{1}{\alpha}\left(1+W_0(e^{-1})\right)\right)\cup\left(-\frac{1}{\alpha}\left(1+W_0(e^{-1})\right),\; \frac{W_0(e^{-1})}{-\alpha}\right)\cup\left(\frac{W_0(e^{-1})}{-\alpha},0\right)$, 
	       then the trivial solution of \eqref{maineq} is asymptotically stable. 
   \item If $\lambda=-\frac{1}{\alpha}\left(1+W_0(e^{-1})\right)$ or $\lambda=0$, then the trivial solution of \eqref{maineq} is stable.
   \item If $\lambda\in(0,\infty)$, then the trivial solution of \eqref{maineq} is unstable.
 \end{enumerate}
\item Consider a large gap size $\beta$ that satisfies $\beta>\frac{\alpha}{W_0(e^{-1})}$. Then, $(1+\beta\lambda)e^{\alpha\lambda}=-1$ at 
 $$ \lambda_{-1}:=-\frac{1}{\beta}+\frac{1}{\alpha}W_{-1}\left(-\frac{\alpha}{\beta}e^{\frac{\alpha}{\beta}}\right) \quad\text{and}\quad 
    \lambda_{0}:=-\frac{1}{\beta}+\frac{1}{\alpha}W_{0}\left(-\frac{\alpha}{\beta}e^{\frac{\alpha}{\beta}}\right), $$ 
		and we have the following subcases.
 \begin{enumerate}
   \item If $\lambda\in\left(-\infty,\; \lambda_{-1}\right)\cup\left(\lambda_{0},\; -\frac{1}{\beta}\right)\cup\left(-\frac{1}{\beta},0\right)$, then the trivial solution 
	       of \eqref{maineq} is asymptotically stable. 
   \item If $\lambda=\lambda_{-1}$, $\lambda=\lambda_{0}$, or $\lambda=0$, then the trivial solution of \eqref{maineq} is stable.
   \item If $\lambda\in\left(\lambda_{-1},\lambda_{0}\right)\cup(0,\infty)$, 
	       then the trivial solution of \eqref{maineq} is unstable.
 \end{enumerate}
\end{enumerate}
\end{theorem}

\begin{remark}
If one uses the nabla backward difference operator on $\p_{\alpha,\beta}$ instead of the Delta forward difference operator, then analogous results may be derived. The nabla differential/difference operator in the nabla case is defined by
\[ x^{\nabla}(t) = \begin{cases} \frac{d}{dt}x(t) &: t\in(k(\alpha+\beta),k(\alpha+\beta)+\alpha] \\ \frac{x(t)-x(t-\beta)}{\beta} &: t=k(\alpha+\beta), \end{cases} \]
and the nabla exponential function is given via
\begin{equation}\label{nablaexpo}
 \widehat{e}_\lambda(t,0) = \frac{e^{\lambda t}}{[(1-\beta\lambda)e^{\beta\lambda}]^k}, \qquad \lambda\in\C\backslash\left\{\frac{1}{\beta}\right\}.
\end{equation}
Thus, for the nabla dynamic equation
\begin{equation}\label{nabeq}
 x^\nabla(t) = \lambda x(t), \quad \lambda\in\C\backslash\left\{\frac{1}{\beta}\right\}, \quad t\in\p_{\alpha,\beta},
\end{equation}
we may compare Theorem \ref{reallambdathm} with the following theorem.
\end{remark}

\begin{equation}\label{complexlam}
 \lambda = -\frac{1}{\beta}+\frac{1}{\alpha}W_z\left(\frac{\rho \alpha}{\beta}e^{\frac{\alpha}{\beta}+i\phi}\right), \quad \phi\in(-\pi,\pi], \quad \rho >0, \quad \beta>0,
\end{equation} 

\begin{theorem}[Delta equation]\label{complexTHM}
Let $(1+\beta\lambda)e^{\alpha\lambda}=\rho e^{i\phi}$, for $\rho >0$ and $\phi\in(-\pi,\pi]$, so that $\lambda\in\C\backslash\left\{-\frac{1}{\beta}\right\}$ is given by \eqref{complexlam}, where we have used the Lambert $W$ function denoted by $W_z$ for any $z\in\Z$. Then, the following hold.
\begin{enumerate}
\item[(a)] If $0<\rho<1$, then the trivial solution of \eqref{maineq} is asymptotically stable.
\item[(b)] If $\rho =1$, then the trivial solution of \eqref{maineq} is stable.
\item[(c)] If $\rho>1$, then the trivial solution of \eqref{maineq} is unstable.
\end{enumerate}
\end{theorem}

\begin{remark}
Consider the nabla case with $\lambda\in\C\backslash\left\{\frac{1}{\beta}\right\}$ on the time scale $\p_{\alpha,\beta}$, for continuous interval size $\alpha>0$ and discrete gap size $\beta>0$. With the nabla exponential function given in \eqref{nablaexpo}, set the base of the exponential function as follows, $(1-\beta\lambda)^{-1}e^{\alpha\lambda}=\rho e^{i\phi}$, for $\rho>0$ and $\phi\in[0,2\pi)$. Again, for various branches of the Lambert $W$ function in the complex plane determined by $z\in\Z$, for $\phi\in[0,2\pi)$, we can solve for $\lambda$ to obtain
\begin{equation}\label{complexnab}
 \lambda = \frac{1}{\beta}-\frac{1}{\alpha}W_z\left(\frac{\alpha}{\rho\beta}e^{\frac{\alpha}{\beta}-i\phi}\right), \quad \phi\in[0,2\pi), \quad \rho >0, \quad \beta>0,
\end{equation} 
with a branch cut along the positive real axis, and principal branch denoted $W_0$. Then,
\begin{equation}\label{thm33eq1}
 \widehat{e}_{\lambda}(t,0)=\left(\frac{e^{\alpha\lambda}}{1-\beta\lambda}\right)^{k}e^{\lambda j}=\left(\rho e^{i\phi}\right)^{k}e^{\lambda j}, \quad j\in[0,\alpha],
\end{equation}
and we may now compare Theorem \ref{complexTHM} with the following theorem.
\end{remark}

\begin{definition}[Hyers--Ulam Stability]
We say that \eqref{nabeq} has Hyers--Ulam stability on $\T$ if and only if given an arbitrary $\varepsilon>0$ and any function $\phi:\T\rightarrow\C$ satisfying 
\begin{equation}\label{phiineq}
 |\phi^\nabla(t)-\lambda\phi(t)|\le\varepsilon,  \quad t\in\T, 
\end{equation}
then there exists a solution $x:\T\rightarrow\C$ of \eqref{nabeq} and a constant $K>0$ such that $|\phi(t)-x(t)|\le K\varepsilon$ for all $t\in\T$. In this case, such a constant $K$ is called an HUS constant for \eqref{nabeq} on $\T$.
\end{definition}

\begin{theorem}\label{complexTHMHUS}
Let $\lambda\in\C\backslash\left\{\frac{1}{\beta}\right\}$ have the form \eqref{complexnab}, and let $W_z$ be the Lambert $W$ function for any $z\in\Z$. 
\begin{enumerate}
\item If $\rho=1$, then \eqref{nabeq} is not Hyers--Ulam stable.
\item If $\rho>1$, then \eqref{nabeq} is Hyers--Ulam stable, with HUS constant at most 
\begin{equation}\label{KRbigger1}
 K_{\rho>1}:=\max_{j\in[0,\alpha]}\left(\frac{\rho-1+e^{j\operatorname{Re}(\lambda)}+Re^{(j-\alpha)\operatorname{Re}(\lambda)}\left(\beta\operatorname{Re}(\lambda)-1\right)}{(\rho-1)\operatorname{Re}(\lambda)}\right), 
\end{equation} 
or $K=\frac{\rho(\beta+\alpha)}{\rho-1}$ if $\operatorname{Re}(\lambda)=0$. 
\item If $0<\rho<1$, then \eqref{nabeq} is Hyers--Ulam stable, with HUS constant at most  
\begin{equation}\label{Krsmaller1}
 K_{0<\rho<1}:= \max_{j\in[0,\alpha]}\left(\frac{\rho-1+e^{j\operatorname{Re}(\lambda)}+Re^{(j-\alpha)\operatorname{Re}(\lambda)}\left(\beta\operatorname{Re}(\lambda)-1\right)}{(1-\rho)\operatorname{Re}(\lambda)}\right). 
\end{equation}
\end{enumerate}
\end{theorem}

\begin{proof}
Let $\lambda\in\C\backslash\left\{\frac{1}{\beta}\right\}$ have the form \eqref{complexnab}, and let $W_z$ be the Lambert $W$ function for any $z\in\Z$. 

Case (i). If $\rho=1$, then $\lambda=\frac{1}{\beta}-\frac{1}{\alpha}W_z\left(\frac{\alpha}{\beta}e^{\frac{\alpha}{\beta}-i\theta}\right)$ for $\theta\in[0,2\pi)$ and for fixed $z\in\Z$. Let the exponential function be given by \eqref{pabexp0}. Then, for $t=k(\alpha+\beta)+j\in[k(\alpha+\beta),k(\alpha+\beta)+\alpha]$ and $j\in[0,\alpha]$, 
\begin{eqnarray*}
 \widehat{e}_{\lambda}(t,0) &=& (1-\beta\lambda)^{-k}e^{k\alpha\lambda} e^{j\lambda} = e^{j\lambda+ik\theta}.
\end{eqnarray*}
Note that, for all $j\in[0,\alpha]$ and $\theta\in[0,2\pi)$, and for any fixed $z\in\Z$, the real part of $\lambda$ satisfies $\operatorname{Re}(\lambda)\ge 0$, and
\[ |\widehat{e}_{\lambda}(t,0) | = e^{j\operatorname{Re}(\lambda)}\in\left[1,\;e^{\alpha\operatorname{Re}(\lambda)}\right]. \]
So, with $\widehat{e}_{\lambda}(t,0)=e^{j\lambda+ik\theta}$ for $t=k(\alpha+\beta)+j$, $j\in[0,\alpha]$, and $\theta\in[0,2\pi)$, set 
\[ \phi(t)=\frac{\varepsilon t\widehat{e}_{\lambda}(t,0)}{e^{\alpha\operatorname{Re}(\lambda)}}. \]
Then, we have
$$ |\phi^{\nabla}(t)-\lambda \phi(t)| = \left|\frac{\varepsilon\lambda t\widehat{e}_{\lambda}(t,0)}{e^{\alpha\operatorname{Re}(\lambda)}}+\frac{\varepsilon e^{\rho}_{\lambda}(t,0)}{e^{\alpha\operatorname{Re}(\lambda)}}-\frac{\varepsilon\lambda t\widehat{e}_{\lambda}(t,0)}{e^{\alpha\operatorname{Re}(\lambda)}}\right| = \frac{\varepsilon}{e^{\alpha\operatorname{Re}(\lambda)}}|e^{\rho}_{\lambda}(t,0)|\le \varepsilon $$
implies that $\phi$ satisfies \eqref{phiineq}, so that
$$ |\phi(t)-x(t)| = \left|\frac{\varepsilon t\widehat{e}_{\lambda}(t,0)}{e^{\alpha\operatorname{Re}(\lambda)}}-x_0\widehat{e}_{\lambda}(t,0)\right| = |\widehat{e}_{\lambda}(t,0)|\left|\frac{\varepsilon t}{e^{\alpha\operatorname{Re}(\lambda)}} - x_0\right|\ge \left|\frac{\varepsilon t}{e^{\alpha\operatorname{Re}(\lambda)}} - x_0\right| \rightarrow\infty $$
for any possible initial condition $x_0$, meaning \eqref{nabeq} is not HUS for $\rho=1$, that is when $\lambda=\frac{1}{\beta}-\frac{1}{\alpha}W_z\left(\frac{\alpha}{\beta}e^{\frac{\alpha}{\beta}-i\theta}\right)$ for any $\theta\in[0,2\pi)$, $\beta>0$, and for any fixed $z\in\Z$.

Case (ii). Let $\rho>1$, that is, let $\lambda = \frac{1}{\beta}-\frac{1}{\alpha}W_z\left(\frac{\alpha}{\beta \rho}e^{\frac{\alpha}{\beta}-i\theta}\right)$, initially with $\operatorname{Re}(\lambda)\ne 0$, for $\theta\in[0,2\pi)$ and $z\in\Z$. Let the exponential function be given by \eqref{pabexp0}, and let $\phi$ satisfy \eqref{phiineq}. Then, $\phi$ has the form given by 
\begin{equation}\label{phieqform}
 \phi(t)=\phi_0\widehat{e}_{\lambda}(t,0)+\widehat{e}_{\lambda}(t,0) \int_0^t\frac{q(s)}{\widehat{e}_{\lambda}(\rho(s),0)}\nabla s, \quad |q(s)|\le\varepsilon\; \forall s\in\T,  
\end{equation}
and again,
$$ \int_0^{\infty} \frac{q(s)}{\widehat{e}_{\lambda}(\rho(s),0)}\nabla s $$
exists and is finite, as $|\widehat{e}_{\lambda}(t,0)|=\rho^ke^{j\operatorname{Re}(\lambda)}$ for $\rho>1$ and $t=k(\alpha+\beta)+j$, $j\in[0,\alpha]$. Note that 
$$ x(t) = x_0\widehat{e}_{\lambda}(t,0), \qquad x_0=\phi_0 + \int_0^{\infty} \frac{q(s)}{\widehat{e}_{\lambda}(\rho(s),0)}\nabla s $$
is a well-defined solution of \eqref{nabeq}.
Now, to integrate from $s=0$ to $s=t=k(\alpha+\beta)+j$ for some $k\in\{0,1,2,\ldots\}$ and $j\in[0,\alpha]$, we see that there are $k$ continuous intervals and $k$ gaps to integrate over, plus the final partial interval (continuous), so that
\begin{eqnarray}
 \int_{0}^{t} \frac{\nabla s}{|\widehat{e}_{\lambda}(\rho(s),0)|} 
  &=& \sum_{m=0}^{k-1}\left(\int_{m(\alpha+\beta)}^{m(\alpha+\beta)+\alpha}\frac{ds}{|\widehat{e}_{\lambda}(s,0)|}\right)+\sum_{m=1}^{k}\left(\int_{m(\alpha+\beta)-\beta}^{m(\alpha+\beta)}\frac{\nabla s}{|\widehat{e}_{\lambda}(\rho(s),0)|}\right)+\int_{k(\alpha+\beta)}^{t}\frac{ds}{|\widehat{e}_{\lambda}(s,0)|} \nonumber\\
	&=& \sum_{m=0}^{k-1}\left(\int_{0}^{\alpha}\frac{e^{-j\operatorname{Re}(\lambda)}}{\rho^{m}}dj\right)+\sum_{m=0}^{k-1}\frac{\beta}{\rho^me^{\alpha\operatorname{Re}(\lambda)}} +\int_{0}^{j}\frac{e^{-\ell\operatorname{Re}(\lambda)}}{\rho^{k}}d\ell \nonumber\\
	&=& \sum_{m=0}^{k-1}\frac{-1}{\rho^m\operatorname{Re}(\lambda)}e^{-j\operatorname{Re}(\lambda)}\bigg|_{j=0}^{\alpha}+\sum_{m=0}^{k-1}\frac{\beta}{\rho^me^{\alpha\operatorname{Re}(\lambda)}}+\frac{-1}{\rho^k\operatorname{Re}(\lambda)}e^{-\ell\operatorname{Re}(\lambda)}\bigg|_{\ell=0}^{j} \nonumber\\
	&=& \frac{\rho(\rho^k-1)(1-e^{-\alpha\operatorname{Re}(\lambda)}))}{\rho^k(\rho-1)\operatorname{Re}(\lambda)} + \frac{e^{-\alpha\operatorname{Re}(\lambda)}\beta \rho(\rho^k-1)}{\rho^k(\rho-1)} + \frac{1-e^{-j\operatorname{Re}(\lambda)}}{\rho^k\operatorname{Re}(\lambda)} \label{int0tot}
\end{eqnarray}
and 
\begin{eqnarray*}
 \int_{0}^{\infty} \frac{\nabla s}{|\widehat{e}_{\lambda}(\rho(s),0)|} 
  = \lim_{t\rightarrow\infty}\int_{0}^{t} \frac{\nabla s}{|\widehat{e}_{\lambda}(\rho(s),0)|} 
	= \frac{\rho(1-e^{-\alpha\operatorname{Re}(\lambda)}))}{(\rho-1)\operatorname{Re}(\lambda)} + \frac{e^{-\alpha\operatorname{Re}(\lambda)}\beta \rho}{\rho-1} .
\end{eqnarray*}
Using these two integral values, we have
\begin{eqnarray*} 
 |\phi(t)-x(t)| &=& |\widehat{e}_{\lambda}(t,0)| \left|-\int_{t}^{\infty} \frac{q(s)}{\widehat{e}_{\lambda}(\rho(s),0)}\nabla s\right| \\
	&\le& \varepsilon |\widehat{e}_{\lambda}(t,0)| \int_{t}^{\infty} \frac{1}{|\widehat{e}_{\lambda}(\rho(s),0)|}\nabla s \\
	& = & \varepsilon |\widehat{e}_{\lambda}(t,0)| \left(\int_{0}^{\infty}-\int_{0}^{t} \right)\frac{1}{|\widehat{e}_{\lambda}(\rho(s),0)|}\nabla s  \\
	& = & \varepsilon \left(\frac{\rho-1+e^{j\operatorname{Re}(\lambda)}+\rho e^{(j-\alpha)\operatorname{Re}(\lambda)}\left(\beta\operatorname{Re}(\lambda)-1\right)}{(\rho-1)\operatorname{Re}(\lambda)}\right)
\end{eqnarray*}
for $j\in[0,\alpha]$, and for fixed $z\in\Z$, $\rho>1$, $\theta\in[0,2\pi)$ that determine $\lambda\in\C$ with $\operatorname{Re}(\lambda)\ne 0$. Set 
$K$ as in \eqref{KRbigger1}, that is,
$$ K_{\rho>1}:=\max_{j\in[0,\alpha]}\left(\frac{\rho-1+e^{j\operatorname{Re}(\lambda)}+\rho e^{(j-\alpha)\operatorname{Re}(\lambda)}\left(\beta\operatorname{Re}(\lambda)-1\right)}{(\rho-1)\operatorname{Re}(\lambda)}\right). $$
Therefore, \eqref{nabeq} has HUS for $\lambda = \frac{1}{\beta}-\frac{1}{\alpha}W_z\left(\frac{\alpha}{\beta \rho}e^{\frac{\alpha}{\beta}-i\theta}\right)$ with $\operatorname{Re}(\lambda)\ne 0$ and $\rho>1$, with HUS constant at most $K_{\rho>1}$. If $\operatorname{Re}(\lambda)=0$, then
\begin{eqnarray*} 
 \lim_{\operatorname{Re}(\lambda)\rightarrow 0} K_{\rho>1} &=& \lim_{\operatorname{Re}(\lambda)\rightarrow 0}\max_{j\in[0,\alpha]}\left(\frac{\rho-1+e^{j\operatorname{Re}(\lambda)}+\rho e^{(j-\alpha)\operatorname{Re}(\lambda)}\left(\beta\operatorname{Re}(\lambda)-1\right)}{(\rho-1)\operatorname{Re}(\lambda)}\right) \\
  &=& \max_{j\in[0,\alpha]}\frac{\rho(\beta+\alpha)+j(1-\rho)}{\rho-1} \\
	&=& \frac{\rho(\beta+\alpha)}{\rho-1},
\end{eqnarray*}
since $\rho>1$ in this case. As a result, \eqref{nabeq} is HUS with HUS constant at most $K_{\rho>1}=\frac{\rho(\beta+\alpha)}{\rho-1}$, for $\rho>1$ and $\operatorname{Re}(\lambda)=0$. In either instance, case (ii) holds.

Case (iii).
Finally, let $\lambda = \frac{1}{\beta}-\frac{1}{\alpha}W_z\left(\frac{\alpha}{\beta \rho}e^{\frac{\alpha}{\beta}-i\theta}\right)$ for $\theta\in[0,2\pi)$ and $\rho\in(0,1)$, let the exponential function be given by \eqref{pabexp0}, and let $\phi$ satisfy \eqref{phiineq}. Using $\rho\in(0,1)$ and \eqref{int0tot}, as well as $t=k(\alpha+\beta)+j$ for $j\in[0,\alpha]$, we can modify \eqref{int0tot} to get
\begin{eqnarray}
 \int_{0}^{t} \frac{\nabla s}{|\widehat{e}_{\lambda}(\rho(s),0)|} 
  &=& \frac{\rho(1-\rho^k)(1-e^{-\alpha\operatorname{Re}(\lambda)})}{\rho^k(1-\rho)\operatorname{Re}(\lambda)} + \frac{e^{-\alpha\operatorname{Re}(\lambda)}\beta \rho(1-\rho^k)}{\rho^k(1-\rho)} + \frac{1-e^{-j\operatorname{Re}(\lambda)}}{\rho^k\operatorname{Re}(\lambda)} \label{int0totsmallR}. 
\end{eqnarray}
If $\phi$ satisfies the perturbed equation \eqref{phiineq}, then $\phi$ is again given as in \eqref{phieqform}.
Let $x$ be a solution of \eqref{nabeq} with form \eqref{xform}, where $x_0=\phi_0-\varepsilon\left(\frac{e^{-\alpha\operatorname{Re}(\lambda)}\beta \rho}{1-\rho}+\frac{\rho(1-e^{-\alpha\operatorname{Re}(\lambda)})}{(1-\rho)\operatorname{Re}(\lambda)}\right)$; note that both fractions in the parentheses here are positive, due to $\rho\in(0,1)$ and $\operatorname{Re}(\lambda)<0$ in this case. Employing \eqref{int0totsmallR} with $t=k(\alpha+\beta)+j$, we see that
\begin{eqnarray*}
 |\phi(t)-x(t)| &=& |\widehat{e}_{\lambda}(t,0)|\left|\phi_0+\int_0^t \frac{q(s)}{\widehat{e}_{\lambda}(\rho(s),0)}\nabla s-\left(\phi_0-\varepsilon\left(\frac{e^{-\alpha\operatorname{Re}(\lambda)}\beta \rho}{1-\rho}+\frac{\rho(1-e^{-\alpha\operatorname{Re}(\lambda)})}{(1-\rho)\operatorname{Re}(\lambda)}\right)\right)\right| \\
 &=& |\widehat{e}_{\lambda}(t,0)|\left|\int_0^t \frac{q(s)}{\widehat{e}_{\lambda}(\rho(s),0)}\nabla s+\varepsilon\left(\frac{e^{-\alpha\operatorname{Re}(\lambda)}\beta \rho}{1-\rho}+\frac{\rho(1-e^{-\alpha\operatorname{Re}(\lambda)})}{(1-\rho)\operatorname{Re}(\lambda)}\right)\right| \\
 &\le& \varepsilon |\widehat{e}_{\lambda}(t,0)|\left(\int_0^t \frac{1}{|\widehat{e}_{\lambda}(\rho(s),0)|}\nabla s + \frac{e^{-\alpha\operatorname{Re}(\lambda)}\beta \rho}{1-\rho}+\frac{\rho(1-e^{-\alpha\operatorname{Re}(\lambda)})}{(1-\rho)\operatorname{Re}(\lambda)}\right) \\
 &=& \varepsilon \rho^ke^{j\operatorname{Re}(\lambda)}\left(\frac{\rho(1-e^{-\alpha\operatorname{Re}(\lambda)})}{\rho^k(1-\rho)\operatorname{Re}(\lambda)} + \frac{e^{-\alpha\operatorname{Re}(\lambda)}\beta \rho}{\rho^k(1-\rho)} + \frac{1-e^{-j\operatorname{Re}(\lambda)}}{\rho^k\operatorname{Re}(\lambda)}\right) \\
 &=& \varepsilon e^{j\operatorname{Re}(\lambda)} \left(\frac{\rho(1-e^{-\alpha\operatorname{Re}(\lambda)})}{(1-\rho)\operatorname{Re}(\lambda)} + \frac{e^{-\alpha\operatorname{Re}(\lambda)}\beta \rho}{1-\rho} + \frac{1-e^{-j\operatorname{Re}(\lambda)}}{\operatorname{Re}(\lambda)}\right) \\
 &\le& \varepsilon\left(\frac{\rho-1+e^{j\operatorname{Re}(\lambda)}+\rho e^{(j-\alpha)\operatorname{Re}(\lambda)}\left(\beta\operatorname{Re}(\lambda)-1\right)}{(1-\rho)\operatorname{Re}(\lambda)}\right)
\end{eqnarray*}
for $j\in[0,\alpha]$, as $\rho\in(0,1)$.
Therefore \eqref{nabeq} has HUS for $\lambda = \frac{1}{\beta}-\frac{1}{\alpha}W_z\left(\frac{\alpha}{\beta \rho}e^{\frac{\alpha}{\beta}-i\theta}\right)$ for $\rho\in(0,1)$, with HUS constant given by at most $K=K_{0<\rho<1}$ given in \eqref{Krsmaller1}.
This ends the proof.
\end{proof}

\end{document}
